\chapter{Burn Excision and Skin Grafting}

\section*{Introduction \& Clinical Indications}
Burn excision and grafting removes nonviable tissue and covers a wound that is unlikely to epithelialize promptly. Deep partial-thickness and full-thickness burns may require excision when spontaneous healing would be prolonged, contracture-prone, or unlikely. Early specialist assessment is essential because burn depth can evolve and operative timing depends on perfusion, resuscitation, infection, available donor skin, and the patient’s overall condition.

Large burns require multidisciplinary care, including airway management, fluid resuscitation, temperature control, analgesia, nutrition, infection prevention, rehabilitation, and psychosocial support. Circumferential eschar may require escharotomy for compromised ventilation or limb perfusion; this is a separate urgent procedure and should not be confused with definitive excision and grafting.

\section*{Relevant Surgical Anatomy}
The surgeon must distinguish viable dermis from leathery eschar while preserving structures that can support graft take. Hands, face, joints, and genitalia have specialized functional requirements. Donor sites are selected according to color, thickness, accessibility, and available uninjured skin. Contracture risk follows relaxed skin tension lines and joint-crossing wounds, so positioning and later therapy influence the reconstruction from the beginning.

\section*{Preoperative Workup \& Preparation}
Assess burn size and depth, airway and inhalation injury, circulation, temperature, urine output, comorbidities, tetanus status, and signs of infection. Photograph and map wounds, plan donor sites, correct hypothermia and major physiologic derangements, and coordinate blood availability for extensive excision. Explain staged operations, pain, scarring, graft loss, infection, donor-site morbidity, contracture, and possible need for dermal substitutes or later reconstruction.

\section*{Surgical Technique \& Procedure}
After preparation, nonviable tissue is excised until a vascular wound bed is reached; tangential excision may preserve viable dermis, whereas deeper excision may be necessary for full-thickness injury. Split-thickness autograft is harvested from an unburned donor site, meshed when expansion is needed, and secured with staples, sutures, or dressings. Hemostasis is meticulous because a bleeding or hematoma-forming bed prevents adherence. Complex wounds may be staged or temporarily covered before definitive autografting.

\section*{Complications \& Risks}
Early risks include blood loss, hypothermia, fluid shifts, infection, graft shear, hematoma, and graft loss. Later problems include hypertrophic scars, pruritus, contracture, pigment change, chronic wounds, and donor-site scarring. Inhalation injury, sepsis, and multiorgan complications may dominate prognosis in extensive burns.

\section*{Postoperative Care \& Recovery}
Protect the graft from shear, maintain appropriate moisture and immobilization, and inspect it according to burn-center protocol. Donor sites require dressings and analgesia. Once safe, early range-of-motion, splinting, pressure therapy, scar management, nutrition, and occupational or physical therapy are critical. Rehabilitation may continue for months or years.

\section*{Outcomes \& Prognosis}
Graft take and long-term function depend on wound-bed quality, burn extent, infection control, nutrition, rehabilitation, and the anatomic site. Early specialist referral and coordinated multidisciplinary care improve the chance of preserving movement and reducing disabling contracture.

\section*{References}
\begin{enumerate}
\item American Burn Association. Guidelines for Burn Care Under Austere Conditions and burn-center referral resources.
\item International Society for Burn Injuries. Practice guidelines for burn care.
\item Herndon DN, ed. Total Burn Care. Elsevier.
\end{enumerate}
